\documentclass[11pt]{article}
\usepackage[margin=0.9in]{geometry}
\usepackage{fontspec}
\usepackage{amsmath,amssymb,amsthm}
\usepackage{xurl}
\begin{document}\sloppy
\begin{center}
{\Large GROK H1315: first six tight $L=10$ cubes for \texttt{a40ff54e580f5d0a} time out; an unbroken enumeration of eight-gate circuits returns no model in 420\,s}\\[4pt]
2026-10-06 13:15 America/Los\_Angeles \quad branch \texttt{main}, tip \texttt{3a74dcb8} \quad automation 41774ea2\\[2pt]
\textbf{P versus NP: OPEN.}
\end{center}

\section*{Peer delta}
Checked at 13:16\,PT. No commit landed after H1258's \texttt{3a74dcb8}; nothing new under \path{research/main/} or \path{research/agents/grok/}. Root \texttt{main.tex} is untouched.

\section*{1. Literature}
None. The next claim concerns this pass's own search.

\section*{2. Pending run read: tight $L=10$ (COMPUTATIONAL, partial)}
\path{l10.out} (H1258, started 12:59:14\,PT, 900\,s cap per cube, 6 workers) at 13:24:42\,PT: cubes 0 to 5 all hit the cap (exit 124); 0 SAT, 0 UNSAT, 6 timeouts out of 75 kept cubes; cubes 6 to 11 running since 13:14\,PT. Timeouts are not counted as UNSAT. So $10\le L(f,2)\le 11$ is unchanged; whether $L(f,2)$ is 10 or 11 is OPEN. At this rate the full pass needs about $75\cdot 900/6$\,s, roughly three hours, before any SPLIT2 refinement.

\section*{3. Smallest honest test: two-pebble schedules over other minimum circuits}
H1812 found 11 steps by the exact two-pebble search (recompute allowed) on one generating eight-gate circuit, and H1948 suggested running the same search over all minimum circuits. A schedule of length 10 on any eight-gate circuit would be a real 10-step program and would give $L(f,2)\le 10$ without the tightness lemma.

\textbf{ENUM8-NULL (COMPUTATIONAL, null).} \path{enum8peb.py} (\texttt{07a5d237}) reuses the \path{exsyn_n.py} clauses at $r=8$ (872 variables, 75239 clauses, recomputed this pass), blocks each found structure, dedupes by the set of gate truth tables, and runs the H1812 BFS with re-simulation on each new circuit. Started 13:17:34\,PT under a 420\,s wall cap; the first CaDiCaL call did not return before the cap (exit 124): 0 circuits enumerated, 0 schedules. The encoding has no symmetry breaking, so even one $r=8$ model is slow here. This gives no information about $L(f,2)$.

\section*{Effect on the standing candidate}
No claim is falsified and no charged resource improves this pass. $10\le L(f,2)\le 11$ stands, the lower bound CONDITIONAL on the tightness lemma (CANDIDATE), the upper bound COMPUTATIONAL (re-simulated at H1258). The largest unconditional confirmed price stays one. The time-state-representation surface stays CANDIDATE. The SAT-to-charged-model compiler and the product-span obstruction are both OPEN.

\section*{Next smallest step}
Seed the circuit enumeration instead of solving cold: assume the H1657 circuit as a starting model, or add topological-order symmetry breaking (non-decreasing first operand), then enumerate and pebble. Keep reading \path{l10.out}; re-simulate any SAT line independently.

\bigskip
\noindent\textbf{Status sentence.} P versus NP remains OPEN: the SAT-to-charged-model compiler and the product-span obstruction are open load-bearing arrows. This pass reads the first six tight $L=10$ cubes for \texttt{a40ff54e580f5d0a} as timeouts (0 SAT, 0 UNSAT) and records that an unbroken enumeration of eight-gate circuits found no model within 420\,s; $10\le L(f,2)\le 11$ is unchanged, with the lower bound CONDITIONAL on the tightness lemma. No claim is falsified, and this is not a proof of anything about P versus NP.
\end{document}
